Comprueba las siguientes afirmaciones sobre familias y parametrizaciones de soluciones. En cada caso, sustituye en la EDO y corrige la afirmación con precisión. La letra $C$ designa una constante real; en el último apartado, $t$ es el parámetro que recorre la curva. No hace falta resolver las EDO desde el principio.
\begin{enumerate}
\item «Para todo $C$, la función $y_C(x)=\tan(x+C)$ es solución de $y'=1+y^2$ en $\mathbb{R}$». Indica los intervalos abiertos máximos de cada miembro.
\item «La familia $y_C(x)=C^2(2+\sin x)$ contiene todas las soluciones de $(2+\sin x)y'=y\cos x$ y representa una sola vez cada dato inicial $y(0)=a$». Decide qué valores de $a$ admite y corrige la parametrización de la constante.
\item «Todas las funciones $y_C(x)=\ln(e^x+C)$ son soluciones de $y'=1-e^{-y}$ en su dominio real». Determina ese dominio según $C$ y averigua qué valores de $C$ hacen verdadera la sustitución.
\item «Al variar $C$, la familia $y_C(x)=Ce^{-x^2}$ recoge todas las soluciones de
\[
y''+4xy'+(4x^2+2)y=0
\]
». Verifica los miembros y contrasta la afirmación con la función $z(x)=xe^{-x^2}$, en $\mathbb{R}$.
\item «La parametrización $x=e^t$, $y=C-e^{-t}$, $t\in\mathbb{R}$, permite representar las soluciones de $x^2y'=1$ tanto para $x>0$ como para $x<0$». Comprueba qué intervalo de $x$ recorre y, si falta una componente del dominio, da una parametrización que la recorra.
\end{enumerate}
